% !TEX TS-program = pdflatex
\documentclass[11pt]{article}

\usepackage[T1]{fontenc}
\usepackage[utf8]{inputenc}
\usepackage{lmodern}
\usepackage{amsmath,amssymb,amsthm,mathtools}
\usepackage{enumitem}
\usepackage{stmaryrd}
\usepackage{xspace}

\newcommand{\Ground}{\mathsf{o}\xspace}
\newcommand{\FV}{\mathsf{FV}\xspace}
\newcommand{\NF}{\mathsf{NF}\xspace}
\newcommand{\Val}{\mathsf{Val}\xspace}
\newcommand{\ObsEq}{\mathrel{\simeq}\xspace}
\newcommand{\ObsLeq}{\mathrel{\lesssim}\xspace}
\newcommand{\hole}[1]{[\,]_{#1}}
\newcommand{\arity}{\operatorname{ar}}
\newcommand{\mctx}{\mathcal C}
\newcommand{\mcell}{\Sigma}
\newcommand{\subst}[2]{[#1/#2]}

\newtheorem{definition}{Definition}
\newtheorem{lemma}{Lemma}
\newtheorem{proposition}{Proposition}
\newtheorem{remark}{Remark}

\title{Contextual Equivalence, Cells, and Semi-Cellular Terms\\
       \large A parametrized formulation and a proof attempt for Lemma 2}
\author{}
\date{}

\begin{document}
\maketitle

\section{Setting}

We work in the simply typed $\lambda$-calculus with one ground type
$\Ground$.  In our case, ground values are
allowed to be first-order terms
\[
        f(a_1,\ldots,a_n)
\]
where $f$ is a first-order symbol of arity $n$ and
$a_1,\ldots,a_n$ are ground constants.  Thus a ground normal form is a
first-order ground term of depth at most 1 (here constants are the special case $n=0$). The set of all such value terms of depth $0$ is written $\Val^0$ and of depth at most $1$ is written $\Val^1$

We assume that all terms are in $\eta$-long form.  As usual,
\[
 \lambda x_1\ldots x_n.u
\]
is an abbreviation for
$\lambda x_1.(\cdots(\lambda x_n.u)\cdots)$, with $u$ of ground type.

\section{Contextual equivalence}

\begin{definition}[Contextual equivalence]
Let $t,t'$ be closed terms of the same type.  We write
\[
        t \ObsLeq^i t'
\]
if, for every closing context $C[-]$ such that $C[t]$ has ground type and
$C[t]\in\Val^i$ we have
\[
        \NF(C[t])=\NF(C[t']).
\]
We let $t\ObsEq t'$ when $t\ObsLeq t'$ and $t'\ObsLeq t$.
Here $\NF(u)$ denotes the $\beta$-normal form of $u$.
\end{definition}

Equivalently, in the usual deterministic simply typed setting, one may
define the relation inductively on types by
\[
\begin{array}{ll}
t,t':\Ground
  &\Longrightarrow\quad
  t\ObsLeq^i t' \iff \NF(t) = \NF(t')\in\Val^i
\\[2mm]
t,t':A\to B
  &\Longrightarrow\quad
  t\ObsLeq t'
  \iff
  \text{for every closed }u:A,\;
  (tu)\ObsLeq^i(t'u).
\end{array}
\]
The first-order ground symbols only affect the possible ground normal
forms; in particular, the observation at ground type is equality of the
resulting first-order ground terms.

\begin{lemma}[Context lemma]
If $u\ObsEq^i u'$, then for every term $t$,
\[
        tu\ObsEq^i tu'.
\]
\end{lemma}

\begin{proposition}[Stretching lemma]
Assume, we work with only one first order constant $f$. Let
$
t = \lambda y_1 \ldots y_n.M[u]
$
be a closed term under observations ${\cal M}$, where $M$ is a context with a hole of ground type.
Then ${\cal M} = {\cal M}_1\cup{\cal M}_2$ and 
\begin{itemize}
\item $t$ is observationally equivalent to
$
\lambda y_1 \ldots y_n.M[M[u]].
$
under ${\cal M}_1$
\item $t$ is observationally equivalent to
$
\lambda y_1 \ldots y_n.fv_1\ldots v_n
$
under ${\cal M}_2$ where $v_i$ is either a constant or $u$ and $\lambda y_1\ldots y_n.u$ is a constant
under observations ${\cal M}_2$. 
\end{itemize}
\end{proposition}
\begin{proof}
Observe that
$
M[M[u]] =  M[x][M[u]/x]
$ 
and we can instantiate the free main variables by $t$ such that $M[u][\vec{t}/\vec{y}]$ is observed. We have
$$
 M[x][M[u]/x][\vec{t}/\vec{y}] = M[\vec{t}/\vec{y}][x][M[u][\vec{t}/\vec{y}]/x]
$$
with 
\begin{itemize}
\item
$
M[\vec{t}/\vec{y}][x] =_\beta x
$
under some observations ${\cal M}'_1$
and
$
M[M[u]][\vec{t}/\vec{y}]
=_\beta
x[M[u][\vec{t}/\vec{y}]/x]
= M[u][\vec{t}/\vec{y}],
$
or
\item
$
M[\vec{t}/\vec{y}][x] =\beta a
$
under some observations ${\cal M}''_1$
and
$
M[M[u]][\vec{t}/\vec{y}]
=_\beta
a[M[u][\vec{t}/\vec{y}]/x]
= a.
$
 or
\item $M[t/y][x] =_\beta fv_1\ldots v_k$ 
under some observations ${\cal M}_2$
where $x\in\FV(v_i)$ or $x\not\in\FV(v_i)$
in the former case, actually $v_i = x$ and then $\lambda y_1\ldots y_n.u$ is a constant under observations ${\cal M}_2$; in the latter case $v_i$ must be a constant since our observable is $f$ applied to constants.
\end{itemize}
The conclusion follows for observations ${\cal M}_1={\cal M}'_1\cup{\cal M}''_1$ and ${\cal M}_2$.
\end{proof}


\section{Parametrized contexts and cells}

We use three parameters:
\[
        m\quad\text{(initial variables)},\qquad
        n\quad\text{(ordinary variables at level 1)},\qquad
        k\quad\text{(ordinary variables at level 0)}.
\]
In addition we operate with an index $K$ that enumerates contextual holes.
The distinction between the three classes is purely in the intent they are used for, but it is useful when imposing explicit bounds on the shape of terms.

Let
\[
 X = \{x_1,\ldots,x_m \},\quad X' = \{x'_1,\ldots,x'_m \}
\]
be the initial variables and
\[
 Y = \{ y_1,\ldots,y_n\},\quad Y' = \{ y'_1,\ldots,y'_n\}
\]
the ordinary variables at level 1 and
\[
 Z = \{ z_1,\ldots,z_k \},\quad Z' = \{ z'_1,\ldots,z'_k \}
\]
the ordinary variables at level 0.
We require that the three sets are disjoint.
The holes
\[
        \hole{1},\ldots,\hole{K}
\]
are always of ground type.

A context may contain variables from the set
\[
 X\cup Y\cup Z
 =
 \{x_1,\ldots,x_m,y_1,\ldots,y_n, z_1,\ldots,z_k\}.
\]

\begin{definition}[$(m,n,k)$-context]
An $(m,n,k)$-context is a term context
\[
        C\hole{1}\ldots\hole{K}
\]
such that
\[
        \FV(C)\subseteq X\cup Y\cup Y'\cup Z\cup Z'
\]
and each hole has ground type.

We write
\[
        \mathcal C_{m,n,k}
\]
for the set of all such contexts.
\end{definition}

The parameter $m$ bounds the number of initial variables, $n$ bounds the
number of ordinary variables of level 1, and $k$ bounds the number of ordinary variables of level 0.

\begin{definition}[$(m,n,K)$-cell]
An $(1,m,n,k)$-cell is a context of the form
\[
\mcell\hole{1}\ldots\hole{K}
 =
 \bigl(
 y\,(C_1\hole{1}\ldots\hole{K})
   \cdots
   (C_p\hole{1}\ldots\hole{K})
 \bigr),
\]
where 
\[
 y\in X\cup Y,
\]
each hole is of ground type, and
\[
 \FV(C_i)=\varnothing.
\]
Moreover, the displayed occurrence of $y$ is the unique head occurrence
of the cell.

We require
\[
        |X|=m,\qquad |Y|=n,\qquad
        |Z|=k.
\]
\end{definition}

In particular, a cell has a unique head variable.  The original definition
of Padovani is recovered by taking $m=0$ and letting the $n$ variables
$y_1,\ldots,y_n$ be the available free variables.  In that formulation a
cell has the form
\[
 \Sigma[\ ]_1\ldots[\ ]_K
 =
 \bigl(y(C_1[\ ]_1\ldots[\ ]_K)\ldots
          (C_p[\ ]_1\ldots[\ ]_K)\bigr),
\]
with the $C_i$ having no free variables other than the variables explicitly
allowed by the surrounding abstraction.

\begin{definition}
    A term $t$ is cellular if and only if it is of the following form:
    \begin{itemize}
    \item $\lambda x_1\ldots x_n.a$ where $a$ is an allowed rhs, or,
    \item $\lambda x_1\ldots x_n.\Sigma[w_1]_1\ldots [w_K]_K$ where:
    \begin{itemize}
    \item $\Sigma$ is a cell whose free variable is amongst $X,Y$,
    \item $w_i = (\lambda\vec{x}\vec{y}\vec{z}.M)x''_1... x''_my''_1...y''_nz''_1...z''_k$
        where $x''_i$ is either $x_i$ or $x'_i$, $y''_i$ is either  $y_i$ or $y'_i$, $z''_i$
        is either $z_i$ or $z'_i$.
        \end{itemize}
        \end{itemize}
\end{definition}



\begin{lemma}
    Every semi-cellular term $t$ is observationally equivalent to a cellular term.
\end{lemma}
\begin{proof}
    By induction on the length of $t$. Assume $t = \lambda x_1\ldots x_n.u$ is in normal
    form, with $u = (x_i u_1\ldots u_p)$. Let $M$ be the minimal context such that
    \begin{itemize}
        \item $M\not= [ ]$,
        \item $u = M [t_1]\ldots[t_K]$
        \item each $t_k$ is a term of the form $\lambda\vec{x}\vec{y}\vec{z}.y\ldots$  with $y\in X\cup Y$.
    \end{itemize}
    For each $k$, $t_k$ is a subterm of some $u_j$, and the closure of this latter term is $k$-cellular. Hence, there exists a cell $\Sigma_k$ and terms $w_{1k},\ldots w_L$
    such that:  
\end{proof}


\end{document}